% Generated draft: requires mathematical and source review.
\documentclass[11pt]{article}
\usepackage[margin=25mm]{geometry}
\usepackage{fontspec}
\setmainfont{Latin Modern Roman}
\usepackage{amsmath,amssymb,mathtools,unicode-math}
\setmathfont{Latin Modern Math}
\usepackage{enumitem,needspace,xurl,hyperref}
\setlength{\parindent}{0pt}
\setlength{\parskip}{5pt}
\setlength{\emergencystretch}{3em}
\newcommand{\N}{\mathbb{N}}
\newcommand{\Z}{\mathbb{Z}}
\newcommand{\Q}{\mathbb{Q}}
\newcommand{\R}{\mathbb{R}}
\newcommand{\C}{\mathbb{C}}
\newcommand{\F}{\mathbb{F}}
\newcommand{\A}{\mathbb{A}}
\newcommand{\PP}{\mathbb{P}}
\newcommand{\E}{\mathbb{E}}
\newcommand{\eps}{\varepsilon}
\newcommand{\dd}{\,\mathrm{d}}
\begin{document}
\section*{UnsolvedMath ID: 11000013}
\textbf{Source problem code:} catalogue entry
\subsection{Definitions and mathematical statement}
Let $S_g$ be a closed orientable genus-$g$ surface with $g\ge3$, let $\operatorname{Mod}_g$ be its orientation-preserving mapping class group, and let
\[
 \mathcal I_g=\ker\bigl(\operatorname{Mod}_g\longrightarrow
 \operatorname{Aut}(H_1(S_g;\Z))\bigr)
\]
be the Torelli group. Two subgroups $A,B$ of a common group are commensurable within that group if $A\cap B$ has finite index in both. Abstract commensurability instead asks only for isomorphic finite-index subgroups; the source does not explicitly choose between these meanings.

\textbf{Question 2.8 (normal subgroups).} If $\Gamma\triangleleft\operatorname{Mod}_g$ is finitely generated, must it be commensurable with $\operatorname{Mod}_g$ or with $\mathcal I_g$? Under the within-group convention, the proposed dichotomy for the intended nontrivial groups is
\[
 [\operatorname{Mod}_g:\Gamma]<\infty
 \quad\text{or}\quad
 [\Gamma:\Gamma\cap\mathcal I_g]<\infty\ \text{ and }\
 [\mathcal I_g:\Gamma\cap\mathcal I_g]<\infty.
\]
A necessary qualification is missing from the literal excerpt: the trivial group is finitely generated and normal but is not commensurable with either of these infinite groups. Thus the substantive question needs a nontriviality or appropriate infinitude restriction; this draft records that issue rather than silently adding it to the source sentence.

The accompanying discussion considers
\[
 D_n=\langle T_c^n:c\text{ an essential simple closed curve on }S_g\rangle
\]
and the normal closure of a single mapping class,
\[
 N_\varphi=\langle h\varphi h^{-1}:h\in\operatorname{Mod}_g\rangle.
\]
It describes these as sources of infinitely generated normal subgroups and asks how to determine the algebraic structure of $N_\varphi$, especially whether it has finite index in $\operatorname{Mod}_g$ or in one of the groups written $\mathcal I_g(k)$. The excerpt supplies neither the restrictions on $n$ or $\varphi$ needed for its infinite-generation description nor the definition of $\mathcal I_g(k)$. In particular, no universal infinite-generation assertion for all powers or all normal closures is adopted. Its final reference to a ``following'' basic test question is a forward reference, not an included test problem.
\subsection{Short English statement}
Apart from the omitted trivial-group exception, must every finitely generated normal subgroup of a genus-at-least-three mapping class group be commensurable with the whole group or the Torelli group? The source also highlights the structure and finite-index behavior of single-element normal closures.
\subsection{Sources}
\begin{itemize}
\item Local UnsolvedMath registry, permanent ID 11000013, source code catalogue entry; no proposer or proposal year supplied.
\item Farb, edited volume \emph{Problems on Mapping Class Groups and Related Topics}, Question 2.8, PDF page 27; \url{https://www.math.uchicago.edu/~farb/papers/mcgbook.pdf}. The registry records PDF-font-structure extraction and access on 2026-07-29. The volume inventory has 329 explicitly labeled Problem, Conjecture, and Question headings, with repeated numbering distinguished by PDF page.
\item The research summary reports advances on normal subgroups since 2006 but provides no individual primary references establishing the status of this finitely generated dichotomy.
\end{itemize}
\end{document}
